\chapter[Cartan--Eilenberg Resolutions and Grothendieck Spectral Sequences]%
  {Cartan--Eilenberg Resolutions\\and Grothendieck Spectral Sequences}
\label{ch:grothendieck}

An injective model computes a total right derived functor. To compute its
cohomology in stages, we resolve each term of a complex while also resolving
its cycles, boundaries, and cohomology. The two filtrations of the resulting
double complex then give distinct first pages. Throughout the right-derived
construction, \(p\) is horizontal cochain degree and \(q\) is vertical
resolution degree. The unsigned horizontal and vertical differentials
commute, and the total differential is \(d_h+(-1)^p d_v\).

\section{Simultaneous resolutions}

For a cochain complex \(K^\bullet\) in an abelian category, write
\[
 B^pK=\operatorname{im}(K^{p-1}\to K^p),\qquad
 Z^pK=\ker(K^p\to K^{p+1}),\qquad
 H^pK=Z^pK/B^pK .
\]
A Cartan--Eilenberg injective resolution is a double complex \(I^{p,q}\)
with \(q\ge0\) and \(p\) bounded below, augmented by a map
\(K^\bullet\to I^{\bullet,0}\) killed by \(d_v\). Its columns
\(I^{p,\bullet}\) and the columns \(B_h^p I\), \(Z_h^p I\), \(H_h^p I\)
obtained by taking horizontal boundaries, cycles, and cohomology are
injective resolutions of \(K^p,B^pK,Z^pK,H^pK\), respectively. Each
augmentation is the one induced by the given map from \(K^\bullet\).
The signed totalization uses the convention of Proposition~\ref{res:B01}.
This definition is \cite[Tag 015H]{stacks-derived}.

\begin{lemma}[Exactness for horizontal invariants]
\label{res:G07}
Suppose \(i:U^\bullet\to V^\bullet\) is injective in each degree and that
\[
      U^p/B^pU\longrightarrow V^p/B^pV
\]
is injective for all \(p\). Put \(W^\bullet=\operatorname{coker}(i)\).
Taking boundaries, cycles, cohomology, or quotients by boundaries of
\(0\to U^\bullet\to V^\bullet\to W^\bullet\to0\) gives a short exact
sequence in each degree. In this chapter an injection satisfying
these two tests is called admissible.
\end{lemma}

\begin{proof}
Identify \(U^\bullet\) with its image and write \(q:V^\bullet\to W^\bullet\)
for the quotient. Generalized elements and lifts along epimorphisms
are used as in the paragraph on subobjects in Chapter~\ref{ch:filtrations}.
A lift can be taken after epic pullback. Image membership tested after
such a pullback implies image membership before it.

The image of \(B^pV\) under \(q\) is \(B^pW\): a representative of
a boundary in \(W^p\) is a differential of a locally lifted
predecessor in \(W^{p-1}\), and that predecessor in turn lifts to
\(V^{p-1}\). The kernel on \(B^pV\) is \(U^p\cap B^pV\). This
intersection equals \(B^pU\) by the injectivity hypothesis on
quotients. Thus boundaries form a short exact sequence.

For the quotient sequence, the first map is monic by hypothesis.
The last map is epic since the composite \(V^p\to W^p\to W^p/B^pW\)
is epic. If a class in \(V^p/B^pV\) maps to zero, lift it to a
representative \(v\) in \(V^p\). Its image \(qv\) belongs to \(B^pW\).
Lift \(qv\) to a boundary \(b\) in \(B^pV\); then \(v-b\) belongs
to \(U^p\) and represents the original class. The composite from
\(U^p/B^pU\) is zero, so this gives precisely the kernel.

For cycles, a lift \(v\in V^p\) of a cycle in \(W^p\) has
\(d_Vv\in U^{p+1}\cap B^{p+1}V=B^{p+1}U\). Lift that boundary to
\(u\in U^p\) with \(d_Uu=d_Vv\). Then \(v-u\) is a cycle lifting
the given one. The kernel on cycles consists of cycles in the
subobject \(U^p\), and they are precisely \(Z^pU\) because the
inclusion in degree \(p+1\) is monic. The first map on cycles is
monic because the inclusion in degree \(p\) is monic.

Finally \(H^pU\to H^pV\) is monic by the already proved quotient
injection, and \(H^pV\to H^pW\) is epic by the cycle lift. If a
cohomology class maps to zero, lift it to a cycle \(v\) in \(V^p\).
Then \(qv\) is a boundary and lifts to a boundary \(b\) in \(V^p\).
The difference \(v-b\) belongs to \(Z^pU\) and gives the original
class. Conversely classes coming from \(U^\bullet\) map to zero.
All maps used are restrictions or induced quotient maps of the
given complex maps, so the short exact sequences have the asserted
natural maps.
\end{proof}

\begin{theorem}[Existence of a Cartan--Eilenberg resolution]
\label{res:G01}
In an abelian category with enough injectives, every bounded-below
cochain complex admits a Cartan--Eilenberg injective resolution.
If \(K^p=0\) for \(p<L\), the resolution can be chosen with
\(I^{p,q}=0\) for \(p<L\).
\end{theorem}

\begin{proof}
Choose injective objects \(E_B^p,E_H^p\) and monomorphisms
\(j_B^p:B^pK\to E_B^p\) and \(j_H^p:H^pK\to E_H^p\).
Choose the target to be zero when the source is zero. Form the complex
\[
   J^p=E_B^p\oplus E_H^p\oplus E_B^{p+1},
   \qquad d_J^p(b,h,b')=(b',0,0).
\]
The first summand is the boundary object and the first two summands
form the cycle object in degree \(p\). By injectivity extend
\(j_B^p\) from \(B^pK\) to a map \(u^p:K^p\to E_B^p\). Also extend
the composite \(Z^pK\twoheadrightarrow H^pK\xrightarrow{j_H^p}E_H^p\)
to a map \(v^p:K^p\to E_H^p\). The third component \(w^p\) is
\(K^p\twoheadrightarrow B^{p+1}K\xrightarrow{j_B^{p+1}}E_B^{p+1}\).
The map \((u^p,v^p,w^p)\) is a chain map. On the image of \(d_K^p\),
the first component in degree \(p+1\) is \(j_B^{p+1}\), while the
second and third vanish. On applying \(d_J^p\), the third component
of the input becomes exactly this first component.

On \(K^p/B^pK\), the last two components give a map to
\(E_H^p\oplus E_B^{p+1}\). On the subobject \(H^pK\) the first of
these components is \(j_H^p\), and on the quotient \(B^{p+1}K\) the
second is \(j_B^{p+1}\). Consequently the map on \(K^p/B^pK\) is
injective by the short exact sequences of these subobjects and
quotients. The map on \(K^p\) is injective by the same argument with
subobject \(B^pK\). We have constructed an admissible injection
\(K^\bullet\to J^\bullet\). It respects both short exact sequences
\[
 0\to B^pK\to Z^pK\to H^pK\to0,\qquad
 0\to Z^pK\to K^p\to B^{p+1}K\to0 .
\]
In \(J\) these become split sequences. This is the simultaneous
injective horseshoe construction for the two sequences, with the
same embedded boundary used on both sides of each differential.

Set \(K_0=K\). Iterate the construction with \(K_q\to J_q\)
and \(K_{q+1}=\operatorname{coker}(K_q\to J_q)\). The vertical
maps of \(I^{p,q}=J_q^p\) are induced by the quotient and embedding
\(J_q\twoheadrightarrow K_{q+1}\hookrightarrow J_{q+1}\). These
are chain maps for \(d_h\), and successive vertical maps compose
to zero. The construction gives short exact sequences with terms
\(K_q,J_q,K_{q+1}\) on horizontal terms and on all three horizontal
invariants by Lemma~\ref{res:G07}. Thus the augmented
vertical complexes are exact, and their terms in nonnegative
vertical degrees are injective finite sums or direct summands of
injectives. This proves all the simultaneous resolution conditions.
For \(p<L\) the first two possible summands of \(J^p\) vanish.
The third vanishes too because \(B^{p+1}K=0\) for \(p+1\le L\).
The same lower bound holds for the quotient at every step.
Compare \cite[Tag 015I]{stacks-derived} and the projective
construction in \cite[Section 5.7]{weibel}.
\end{proof}

\section{Comparison and totalization}

\begin{lemma}[Extension into a split row]
\label{res:G08}
A horizontal row of a Cartan--Eilenberg injective resolution has the
extension property along admissible injections of complexes.
Moreover, in the augmented vertical resolution, the inclusion of
each vertical kernel into its row is admissible.
\end{lemma}

\begin{proof}
In a fixed row, the short exact sequences for cycles and boundaries
split because the subobjects \(B_h^p I\) and \(Z_h^p I\) are injective.
With a choice of splitting the row has the form used in
Theorem~\ref{res:G01}. It is the degreewise product of stalks on
injective cohomology objects and two-term complexes on injective
boundary objects, with identity as the only nonzero differential in
each two-term complex. Only three factors can contribute in a given
degree, so the product exists in the category of complexes using
only finite products in the abelian category. A map to a stalk of
an injective \(E\) in degree \(p\) is a map \(U^p/B^pU\to E\). A
map to the two-term identity complex in degrees \(p-1,p\) is
determined by its degree \(p\) component \(U^p\to E\), with its
degree \(p-1\) component obtained by precomposing with \(d_U^{p-1}\).
Both components extend along the corresponding monomorphisms into
\(V\). Taking the product of these extensions proves the claim.

For a vertical degree \(q\ge0\), put
\(T_q=\ker(d_v:I^{\bullet,q}\to I^{\bullet,q+1})\), so \(T_q\)
is also the previous vertical image. At \(q=0\) this uses the
augmentation from \(K\). Exactness of the augmented boundary
resolutions shows that
\[
  T_q^p\cap B_h^p(I^{\bullet,q})
    = \operatorname{im}\bigl(d_v:B_h^p(I^{\bullet,q-1})
                           \to B_h^p(I^{\bullet,q})\bigr)
    = B^p T_q ,
\]
with \(B^pK\) and its augmentation in place of the previous row at
\(q=0\). For the last equality, take the horizontal boundary of
the image of the previous vertical map and use that its
horizontal differential is induced from the previous row.
This equality of subobjects says exactly that
\(T_q^p/B^pT_q\to I^{p,q}/B_h^p(I^{\bullet,q})\) is injective.
The inclusion is thus admissible.
\end{proof}

\begin{theorem}[Comparison and injective totalization]
\label{res:G02}
Let \(I\) and \(J\) be Cartan--Eilenberg resolutions of bounded-below
complexes \(K\) and \(L\). Every map \(f:K\to L\) lifts to a map
of double complexes \(I\to J\). Any two lifts of \(f\) are
vertically homotopic through maps of horizontal complexes.
The comparison of total complexes is compatible with the
augmentations, which are quasi-isomorphisms. Each total complex is
bounded below and termwise injective. For an additive functor \(U\)
between abelian categories, there is a natural chain isomorphism
\[
        U(\operatorname{Tot}I)\simeq\operatorname{Tot}(U I)
\]
preserving the two split filtrations from the diagonal sums.
\end{theorem}

\begin{proof}
First extend the composite \(K\to L\to J^{\bullet,0}\)
along \(K\hookrightarrow I^{\bullet,0}\) using
Lemma~\ref{res:G08}. Given the comparison
on row \(q\) preserving the preceding vertical image, it induces
a map on the quotient of that row by its vertical kernel. The
quotient is the next vertical kernel. Compose its induced map
with the inclusion into the next row of \(J\) and extend along
the admissible inclusion into the next row of \(I\). This constructs
the lift with all vertical squares commuting. Each extension is
a map of horizontal complexes.

Let \(a,b:I\to J\) be two lifts. We construct horizontal chain
maps \(s^q:I^{\bullet,q}\to J^{\bullet,q-1}\) with \(s^0=0\).
In degree zero, \(a^0-b^0\) vanishes on \(K\), so it factors
through the first vertical image of \(I\). Extend this factor to
the whole first row with target \(J^{\bullet,0}\) to obtain \(s^1\).
At the next step consider
\[
      a^q-b^q-d_v^J s^q .
\]
Its composite with \(d_v^I\) from row \(q-1\) is zero. Indeed,
commutation with \(d_v\) and the preceding step make this
composite \(d_v^Jd_v^J s^{q-1}\). The residual map therefore
factors through the next vertical image, and extends to a map
\(s^{q+1}\) by the extension property. This gives
\(a-b=d_v^J s+s\,d_v^I\). On total complexes the homotopy sends
the component of bidegree \((p,q)\) to \((-1)^p s^q\).
In \(D_J((-1)^p s)+((-1)^p s)D_I\) the horizontal terms cancel
since \(s\) commutes with \(d_h\) and a horizontal step changes
the sign. The two vertical terms have the product of two
factors \((-1)^p\), leaving \(d_v^J s+s\,d_v^I\).

Every vertical resolution has cohomology only in degree zero
with the given augmentation from \(K^p\) or \(L^p\). Its total
augmentation is a quasi-isomorphism by Theorem~\ref{res:B04}.
In degree \(n\) the total complex has summands only for
\(L_0\le p\le n\) and \(q=n-p\), for a fixed lower bound \(L_0\).
Thus finite direct sums of injectives give degreewise injectivity.
For two resolutions of the same complex, lifts of the identity
in both directions have composites vertically homotopic to the
identities by the preceding uniqueness proof. Finally an additive
functor preserves finite biproducts and signs. The canonical
biproduct isomorphism commutes with both component differentials
and the inclusions defining the row and column filtrations. This
proves the assertion about \(U\). See \cite[Tags 015J and 015K]{stacks-derived}.
\end{proof}

\section{The two hypercohomology spectral sequences}

\begin{theorem}[Two filtrations for a total right derived functor]
\label{res:G03}
Let \(F:\mathcal A\to\mathcal B\) be left exact between abelian
categories, with \(\mathcal A\) having enough injectives. For a bounded-below
complex \(K^\bullet\), there are spectral sequences with
\[
 {}_{\mathrm c}E_1^{p,q}=R^qF(K^p),
       \qquad
 {}_{\mathrm r}E_2^{p,q}=R^pF(H^qK).
\]
Both converge to \(H^{p+q}(RF(K))\) with finite filtrations in
each target degree. On the first page of the column sequence,
\(d_1\) is induced by \(d_K\). The column sequence is natural in
chain maps from \(E_1\) onward. The row sequence is natural in
chain maps from \(E_2\) onward and descends from that page,
together with its filtered target, to \(D^+(\mathcal A)\).
\end{theorem}

\begin{proof}
Choose a Cartan--Eilenberg resolution \(K\to I\). The
quasi-isomorphism \(K\to\operatorname{Tot}I\) is a bounded-below
injective replacement by Theorem~\ref{res:G02}. Thus
\(\operatorname{Tot}(FI)\simeq F(\operatorname{Tot}I)\) computes
\(RF(K)\) by Theorem~\ref{res:R06}. The double complex has finite
diagonals and the same lower bound. Apply Theorems~\ref{res:S01}
and~\ref{res:B04}, and transport the filtrations on total
cohomology via this comparison.

In the column filtration, Proposition~\ref{res:B03} identifies
the first page in bidegree \((p,q)\) with the cohomology of
\(F(I^{p,\bullet})\) in degree \(q\). The augmented column is an
injective resolution of \(K^p\), so this is \(R^qF(K^p)\) by
Theorem~\ref{res:R09}. The unsigned horizontal differential
is a chain map of these resolutions lifting \(d_K^p\), so the
induced first differential is \(R^qF(d_K^p)\).

For the row filtration, in every fixed vertical degree the
horizontal row is split into stalks and identity pairs as in
Lemma~\ref{res:G08}. Additivity of \(F\)
preserves the two splitting sequences and their differentials.
Therefore
\[
 {}_{\mathrm r}E_1^{p,q}
    \simeq H_h^q(FI^{\bullet,p})
    \simeq F(H_h^q I^{\bullet,p}).
\]
The second identification can be defined by the maps on the
split kernel, image, and quotient. It is natural in horizontal
chain maps and independent of a splitting. In the signed
interchange coordinates of Proposition~\ref{res:B03}, the
first differential is \(F\) of the induced vertical map on
\(H_h^q I^{\bullet,p}\). In the original coordinates it has the
additional factor \((-1)^q\). The vertical cohomology of
\(F(H_h^q I^{\bullet,*})\) computes \(R^pF(H^qK)\) by the
simultaneous resolution condition. This gives the row second page.

For a chain map, take a lift between the double resolutions.
It preserves each filtration because it preserves both indices.
Two lifts differ by a vertical homotopy commuting with the
horizontal differentials. In each column this is a homotopy
after applying \(F\), so the induced column map on \(E_1\) is
independent of the lift. On the horizontal cohomology of each
row it gives a homotopy in the vertical direction and hence the
same row map on \(E_2\). The maps on later pages are induced
iteratively by homology. The total homotopy gives equality on
underlying cohomology. Both lifted maps preserve the filtrations,
so their restrictions to the filtration steps coincide as well.
Lifting identities and composites then proves naturality.
In particular, different choices of the double resolution
are compared by lifts of the identity, inverse on the pages
in question and on both filtered targets.

For the row sequence a quasi-isomorphism of bounded-below complexes
induces isomorphisms on \(H^q K\) and therefore on \(E_2\).
Finite convergence and Corollary~\ref{res:S03} give an isomorphism
on the filtered target. Thus the row construction inverts
quasi-isomorphisms on the stated pages and admits a unique
descended functor by the localization universal property.
For compatibility with homotopy classes, use the cylinder
\(K\oplus\operatorname{Cone}(\mathrm{id}_K)\), whose projection to
\(K\) is a quasi-isomorphism by Proposition~\ref{res:K03}. With cone
coordinates \((a,b)\) and differential \((da+b,-db)\), its two
inclusions send \(x\) to \((x,(0,0))\) and \((x,(x,0))\).
They induce the same row maps because each is a section of
the projection. If \(g-f=dh+hd\), the map
\((x,(a,b))\mapsto f(x)+(g-f)(a)+h(b)\) is a chain map to the
common target of \(f,g\). The homotopy identity checks the
chain map property on the second cone coordinate. Its
restrictions are \(f\) and \(g\), proving compatibility.
The column first page uses the terms of \(K\) and its
functoriality statement is for complexes. These are the two computations in
\cite[Tags 015J and 015K]{stacks-derived}.
\end{proof}

\section{A composition of right derived functors}

\begin{theorem}[Grothendieck spectral sequence]
\label{res:G04}
Let \(F:\mathcal A\to\mathcal B\) and \(G:\mathcal B\to\mathcal C\)
be left exact between abelian categories. Suppose \(\mathcal A\) and
\(\mathcal B\) have enough injectives and \(F(I)\) is \(G\)-acyclic
for every injective \(I\) of \(\mathcal A\). For \(X\in D^+(\mathcal A)\)
there is a natural cohomological spectral sequence
\[
      E_2^{p,q}=R^pG\bigl(H^q(RF(X))\bigr)
      \quad\Longrightarrow\quad
      H^{p+q}\bigl(R(G\circ F)(X)\bigr)
\]
with differential of bidegree \((r,1-r)\) on page \(r\).
The target filtration is finite in each degree.
\end{theorem}

\begin{proof}
Choose a bounded-below injective model \(I^\bullet\) of \(X\).
Then \(F(I^\bullet)\) is a bounded-below representative of \(RF(X)\).
The acyclicity hypothesis means that every term of this representative
computes its own derived \(G\): its higher \(R^iG\) vanish, and
the degree-zero comparison is the isomorphism given by left exactness.
The acyclic-complex comparison of Theorem~\ref{res:R08} identifies
\(G(F(I^\bullet))\) with \(RG(F(I^\bullet))\) in the derived category.
This is the component of the canonical composition map of
Theorem~\ref{res:R12},
\[
           R(G\circ F)(X)\longrightarrow RG(RF(X)).
\]
It is an isomorphism natural in \(X\).

Apply the row sequence of Theorem~\ref{res:G03} to the functor \(G\)
and the bounded-below complex \(F(I^\bullet)\). Its second page is
the one displayed above, and it converges to \(H^*(RG(RF(X)))\).
Pull back its finite filtration along the canonical composition
isomorphism. Since the row construction from \(E_2\) onward is
functorial on the derived category, the choice of \(I^\bullet\)
or of a resolution for \(F(I^\bullet)\) gives the same comparisons
from the stated page onward. The filtration is finite because
the resolution degree is nonnegative and the horizontal degree
has a fixed lower bound. See \cite[Tags 015M and 015N]{stacks-derived}
and \cite[Section 5.8]{weibel}.
\end{proof}

\begin{corollary}[Object case, edge maps, and low degrees]
\label{res:G05}
Under the hypotheses of Theorem~\ref{res:G04}, an object \(A\) of
\(\mathcal A\), placed in degree zero, gives a first-quadrant sequence
\[
      E_2^{p,q}=R^pG(R^qF(A))
       \quad\Longrightarrow\quad R^{p+q}(G\circ F)(A).
\]
The natural edge maps in total degree \(n\ge0\) have the form
\[
 R^nG(F(A))\ \longrightarrow\ R^n(G\circ F)(A)\
                    \longrightarrow\ G(R^nF(A)).
\]
In degrees one and two they occur in the exact sequence
\[
\begin{split}
 0\longrightarrow R^1G(F(A))
 &\longrightarrow R^1(G\circ F)(A)
 \longrightarrow G(R^1F(A))\\
 &\xrightarrow{d_2^{0,1}} R^2G(F(A))
 \longrightarrow R^2(G\circ F)(A).
\end{split}
\]
\end{corollary}

\begin{proof}
Write \(Y=RF(A[0])\). It has cohomology in nonnegative degrees,
with \(H^0Y=F(A)\) and \(H^qY=R^qF(A)\) by
Theorem~\ref{res:R09}. This gives the quadrant assertion.
The first edge map is induced by the canonical map
\[
       F(A)[0]\simeq\tau_{\le0}Y\longrightarrow Y
\]
after applying \(RG\), taking \(H^n\), and identifying the
target via Theorem~\ref{res:G04}. The row sequence on \(F(A)[0]\)
has only its \(q=0\) row at \(E_2\). Its map to the sequence for
\(Y\) is the identity on the second-page term in bidegree
\((n,0)\). On cohomology it factors through the highest possibly
nonzero filtration step \(F^nH^n\), identifying the map as
\(E_2^{n,0}\twoheadrightarrow E_\infty^{n,0}
  \simeq F^nH^n\hookrightarrow H^n\).

For the second edge map use
\[
      Y\ \longleftarrow\ \tau_{\le n}Y\
               \longrightarrow H^n(Y)[-n].
\]
The left arrow induces an isomorphism on \(H^nRG\): its
truncation triangle has third term \(\tau_{\ge n+1}Y\), and
choosing a bounded-below injective model in degrees at least
\(n+1\) by Theorem~\ref{res:R01} shows its \(RG\) has zero cohomology
in degrees \(n-1,n\). Take the inverse isomorphism on
\(H^nRG\) and then the map to
\(H^nRG(H^n(Y)[-n])=G(H^n(Y))\).
By row functoriality the two arrows are the identity on the
\(E_2^{0,n}\) term, and the sequence for \(H^n(Y)[-n]\) has
only its \(q=n\) row at \(E_2\). In total degree \(n\) its
only term has filtration index zero. Thus the resulting map
is the edge map
\(H^n\twoheadrightarrow E_\infty^{0,n}\hookrightarrow E_2^{0,n}\).
This description also proves naturality of both edges.

In total degree one the graded pieces are
\(E_\infty^{1,0}=E_2^{1,0}\) and
\(E_\infty^{0,1}=\ker(d_2^{0,1})\). The only differential
that can enter \((2,0)\) from page two onward is \(d_2^{0,1}\);
none can leave \((2,0)\). Hence \(E_\infty^{2,0}\) is the
cokernel of this differential and injects into total degree two.
These identifications give exactness at all intervening terms
and no terminal surjectivity assertion. This is the filtration
argument of Theorem~\ref{res:S05} in the present sequence.
See \cite[Tag 015N]{stacks-derived} and
\cite[Theorem 5.8.3]{weibel}.
\end{proof}

\section{The homological dual}

\begin{theorem}[Grothendieck spectral sequence for left derivation]
\label{res:G06}
Let \(F:\mathcal A\to\mathcal B\) and \(G:\mathcal B\to\mathcal C\)
be right exact between abelian categories. Suppose \(\mathcal A\) and
\(\mathcal B\) have enough projectives and \(L_iG(F(P))=0\) for
\(i>0\) and projective \(P\) of \(\mathcal A\). For an object \(A\)
there is a natural first-quadrant homological spectral sequence
\[
      E^2_{p,q}=L_pG(L_qF(A))
       \quad\Longrightarrow\quad L_{p+q}(G\circ F)(A).
\]
More generally, for a derived object represented by a bounded-above
cochain complex \(X\), use homological grading \(H_j=H^{-j}\).
The second page is \(L_pG(H_q(LF(X)))\) and the target is
\(H_{p+q}(L(G\circ F)(X))\). The differential has bidegree
\((-r,r-1)\) and the target filtration is increasing and finite
in each degree.
\end{theorem}

\begin{proof}
Pass to the opposite abelian categories. The opposite functors
are left exact and the opposite source and intermediate
categories have enough injectives. Start with the chain indices
\(K_j=X^{-j}\), with differential \(\partial_j:K_j\to K_{j-1}\).
Viewing \(K_j\) in the opposite category gives the degree \(j\)
term of a cochain complex with differential \(\partial_{j+1}^{\mathrm{op}}\).
The chain lower bound becomes a cochain lower bound there.
Reversing arrows interchanges kernels with cokernels and
carries the degree \(j\) chain homology to degree \(j\)
cohomology in the opposite category. Projective resolutions
become injective resolutions. After reversing back, these
are the bounded-above projective models of
Theorem~\ref{res:R07}. The right acyclicity hypothesis of
Theorem~\ref{res:G04} is exactly the displayed vanishing,
with the degree-zero comparison supplied by right exactness.

Under arrow reversal the canonical comparison becomes
\[
              LG(LF(X))\longrightarrow L(G\circ F)(X).
\]
On a bounded-above projective model it is the canonical
derived comparison for applying \(G\) to a termwise
\(G\)-acyclic complex, hence an isomorphism also by
Theorem~\ref{res:R08}. Applying Theorem~\ref{res:G04}
in the opposite categories and reversing the arrows gives
the specified pages and target. A cohomological differential
from bidegree \((p,q)\) to \((p+r,q-r+1)\) becomes a homological
differential from the latter bidegree to the former. Equivalently,
the resulting page differential has bidegree \((-r,r-1)\).
The signed total in the opposite categories dualizes termwise
to \(d_h+(-1)^p d_v\) with both indices lowered. This agrees with
the chain convention of Chapter~\ref{ch:double-complexes}.

For the filtration, a decreasing subobject \(F^p\) of a target
in the opposite category becomes a quotient of the target in
the original category. Define the increasing step of index \(p-1\)
to be the kernel of that quotient. Reversing the short exact
sequence of successive decreasing steps identifies the
quotient of the increasing steps of indices \(p,p-1\) with the
reversed \(p\)-th graded piece. Finiteness follows from
finiteness in the opposite category. The object \(A[0]\) has
the nonnegative resolution degrees in both axes, giving the
first quadrant and the claimed object formula.

In particular, for \(A[0]\) the right-derived edge maps of
Corollary~\ref{res:G05} reverse to
\[
 G(L_nF(A))\longrightarrow L_n(G\circ F)(A)
                     \longrightarrow L_nG(F(A)).
\]
The five-term sequence reverses to
\[
\begin{split}
 L_2(G\circ F)(A)\longrightarrow L_2G(F(A))
 &\xrightarrow{d_2}G(L_1F(A))
 \longrightarrow L_1(G\circ F)(A)\\
 &\longrightarrow L_1G(F(A))\longrightarrow0 .
\end{split}
\]
Here \(d_2\) goes from \((2,0)\) to \((0,1)\). This is the
homological case of \cite[Corollary 5.8.4]{weibel}.
\end{proof}
